% Options for packages loaded elsewhere
% Options for packages loaded elsewhere
\PassOptionsToPackage{unicode}{hyperref}
\PassOptionsToPackage{hyphens}{url}
\PassOptionsToPackage{dvipsnames,svgnames,x11names}{xcolor}
%
\documentclass[
  letterpaper,
  DIV=11,
  numbers=noendperiod]{scrartcl}
\usepackage{xcolor}
\usepackage{amsmath,amssymb}
\setcounter{secnumdepth}{-\maxdimen} % remove section numbering
\usepackage{iftex}
\ifPDFTeX
  \usepackage[T1]{fontenc}
  \usepackage[utf8]{inputenc}
  \usepackage{textcomp} % provide euro and other symbols
\else % if luatex or xetex
  \usepackage{unicode-math} % this also loads fontspec
  \defaultfontfeatures{Scale=MatchLowercase}
  \defaultfontfeatures[\rmfamily]{Ligatures=TeX,Scale=1}
\fi
\usepackage{lmodern}
\ifPDFTeX\else
  % xetex/luatex font selection
\fi
% Use upquote if available, for straight quotes in verbatim environments
\IfFileExists{upquote.sty}{\usepackage{upquote}}{}
\IfFileExists{microtype.sty}{% use microtype if available
  \usepackage[]{microtype}
  \UseMicrotypeSet[protrusion]{basicmath} % disable protrusion for tt fonts
}{}
\makeatletter
\@ifundefined{KOMAClassName}{% if non-KOMA class
  \IfFileExists{parskip.sty}{%
    \usepackage{parskip}
  }{% else
    \setlength{\parindent}{0pt}
    \setlength{\parskip}{6pt plus 2pt minus 1pt}}
}{% if KOMA class
  \KOMAoptions{parskip=half}}
\makeatother
% Make \paragraph and \subparagraph free-standing
\makeatletter
\ifx\paragraph\undefined\else
  \let\oldparagraph\paragraph
  \renewcommand{\paragraph}{
    \@ifstar
      \xxxParagraphStar
      \xxxParagraphNoStar
  }
  \newcommand{\xxxParagraphStar}[1]{\oldparagraph*{#1}\mbox{}}
  \newcommand{\xxxParagraphNoStar}[1]{\oldparagraph{#1}\mbox{}}
\fi
\ifx\subparagraph\undefined\else
  \let\oldsubparagraph\subparagraph
  \renewcommand{\subparagraph}{
    \@ifstar
      \xxxSubParagraphStar
      \xxxSubParagraphNoStar
  }
  \newcommand{\xxxSubParagraphStar}[1]{\oldsubparagraph*{#1}\mbox{}}
  \newcommand{\xxxSubParagraphNoStar}[1]{\oldsubparagraph{#1}\mbox{}}
\fi
\makeatother


\usepackage{longtable,booktabs,array}
\newcounter{none} % for unnumbered tables
\usepackage{calc} % for calculating minipage widths
% Correct order of tables after \paragraph or \subparagraph
\usepackage{etoolbox}
\makeatletter
\patchcmd\longtable{\par}{\if@noskipsec\mbox{}\fi\par}{}{}
\makeatother
% Allow footnotes in longtable head/foot
\IfFileExists{footnotehyper.sty}{\usepackage{footnotehyper}}{\usepackage{footnote}}
\makesavenoteenv{longtable}
\usepackage{graphicx}
\makeatletter
\newsavebox\pandoc@box
\newcommand*\pandocbounded[1]{% scales image to fit in text height/width
  \sbox\pandoc@box{#1}%
  \Gscale@div\@tempa{\textheight}{\dimexpr\ht\pandoc@box+\dp\pandoc@box\relax}%
  \Gscale@div\@tempb{\linewidth}{\wd\pandoc@box}%
  \ifdim\@tempb\p@<\@tempa\p@\let\@tempa\@tempb\fi% select the smaller of both
  \ifdim\@tempa\p@<\p@\scalebox{\@tempa}{\usebox\pandoc@box}%
  \else\usebox{\pandoc@box}%
  \fi%
}
% Set default figure placement to htbp
\def\fps@figure{htbp}
\makeatother





\setlength{\emergencystretch}{3em} % prevent overfull lines

\providecommand{\tightlist}{%
  \setlength{\itemsep}{0pt}\setlength{\parskip}{0pt}}



 


\KOMAoption{captions}{tableheading}
\makeatletter
\@ifpackageloaded{caption}{}{\usepackage{caption}}
\AtBeginDocument{%
\ifdefined\contentsname
  \renewcommand*\contentsname{Table of contents}
\else
  \newcommand\contentsname{Table of contents}
\fi
\ifdefined\listfigurename
  \renewcommand*\listfigurename{List of Figures}
\else
  \newcommand\listfigurename{List of Figures}
\fi
\ifdefined\listtablename
  \renewcommand*\listtablename{List of Tables}
\else
  \newcommand\listtablename{List of Tables}
\fi
\ifdefined\figurename
  \renewcommand*\figurename{Figure}
\else
  \newcommand\figurename{Figure}
\fi
\ifdefined\tablename
  \renewcommand*\tablename{Table}
\else
  \newcommand\tablename{Table}
\fi
}
\@ifpackageloaded{float}{}{\usepackage{float}}
\floatstyle{ruled}
\@ifundefined{c@chapter}{\newfloat{codelisting}{h}{lop}}{\newfloat{codelisting}{h}{lop}[chapter]}
\floatname{codelisting}{Listing}
\newcommand*\listoflistings{\listof{codelisting}{List of Listings}}
\makeatother
\makeatletter
\makeatother
\makeatletter
\@ifpackageloaded{caption}{}{\usepackage{caption}}
\@ifpackageloaded{subcaption}{}{\usepackage{subcaption}}
\makeatother
\usepackage{bookmark}
\IfFileExists{xurl.sty}{\usepackage{xurl}}{} % add URL line breaks if available
\urlstyle{same}
% fallback for those not using the hyperref driver hyperxmp:
\makeatletter
\@ifundefined{xmpquote}{\newcommand{\xmpquote}[1]{#1}}{}
\makeatother
\hypersetup{
  pdftitle={2.1 Intuición geométrica: el determinante como factor de escala},
  colorlinks=true,
  linkcolor={blue},
  filecolor={Maroon},
  citecolor={Blue},
  urlcolor={Blue},
  pdfcreator={LaTeX via pandoc}}


\title{2.1 Intuición geométrica: el determinante como factor de escala}
\author{}
\date{}
\begin{document}
\maketitle


\subsection{2.1.1 ¿Qué le hace una matriz al
espacio?}\label{quuxe9-le-hace-una-matriz-al-espacio}

En la Unidad 1 aprendimos a resolver sistemas de ecuaciones y a operar
con matrices. Pero hay una pregunta que casi no hicimos: \textbf{¿qué le
hace una matriz al espacio geométrico?}

Imagina que cada punto \((x, y)\) del plano es un vector. Si multiplicas
todos los puntos del plano por una matriz
\(A = \begin{pmatrix} a & b \\ c & d \end{pmatrix}\), el plano se
estira, se rota, se inclina o se aplasta. La matriz es una
\textbf{transformación lineal}.

Para medir cuánto ``deforma'' la matriz el espacio, necesitamos un
número que resuma ese efecto. Ese número es el \textbf{determinante}.
Antes de aprender a calcularlo con fórmulas, vamos a descubrirlo
mirando.

\subsection{\texorpdfstring{2.1.2 El caso \(2 \times 2\): Áreas con
signo}{2.1.2 El caso 2 \textbackslash times 2: Áreas con signo}}\label{el-caso-2-times-2-uxe1reas-con-signo}

Toma la matriz identidad
\(I = \begin{pmatrix} 1 & 0 \\ 0 & 1 \end{pmatrix}\). Si multiplicas
cualquier vector por \(I\), el vector no cambia. El espacio se queda
exactamente igual. Queremos que el determinante de la identidad sea
\(1\), porque el ``factor de escala'' del espacio es \(1\) (no hay
escala).

Ahora, mira las columnas de una matriz
\(A = \begin{pmatrix} a & b \\ c & d \end{pmatrix}\) como dos vectores:
\(\mathbf{u} = (a, c)\) y \(\mathbf{v} = (b, d)\). Estos dos vectores,
saliendo desde el origen, forman un \textbf{paralelogramo}.

\textbf{Definición 2.1.1 (Determinante de \(2 \times 2\)).} El
determinante de la matriz \(A\) se define como:
\[ \det(A) = \begin{vmatrix} a & b \\ c & d \end{vmatrix} = ad - bc \]

\textbf{Teorema 2.1.1 (Interpretación geométrica).} El valor absoluto
del determinante de \(A\) es exactamente el \textbf{área del
paralelogramo} formado por sus columnas \(\mathbf{u}\) y \(\mathbf{v}\).
\[ \text{Área} = |\det(A)| = |ad - bc| \]

Veamos un ejemplo. Sea
\(A = \begin{pmatrix} 3 & 1 \\ 1 & 2 \end{pmatrix}\). Sus columnas son
\(\mathbf{u} = (3, 1)\) y \(\mathbf{v} = (1, 2)\).
\[ \det(A) = (3)(2) - (1)(1) = 6 - 1 = 5 \] El paralelogramo formado por
estos vectores tiene un área de exactamente \(5\) unidades cuadradas.

\begin{figure}

\centering{

\pandocbounded{\includegraphics[keepaspectratio]{2.1-intuicion-geometrica_files/figure-pdf/fig-paralelogramo-2d-output-1.pdf}}

}

\caption{\label{fig-paralelogramo-2d}El determinante como área. Las
columnas de la matriz forman un paralelogramo. Su área es \(|ad - bc|\).
Aquí, el área es \(5\).}

\end{figure}%




\end{document}
